% =============================================================================
% Chapter 6  Systems and Nonlinearity  (6_extensions.tex)  --  ~6 pp
% rubric: Project Body (Modelling) -- Model Formulation /30.
% Plan: ThesisPlanning.md §5.6 v3.2 (T-31, T-32; supporting contribution S4).
%
% PURPOSE     Take the construction from one real scalar field, Wu et al.'s
%             scope, to several coupled fields, to complex-valued linear
%             equations, and to polynomial nonlinearity, and answer Q(iii):
%             separate what the doubled space can claim (an exact encoding, a
%             projection) from what the Carleman lift can (a preparation, in
%             both forms), and set up the fluids. Q(iii)'s second half (does the
%             reformulation's advantage survive there?; StructuralReview.md F5) is
%             answered in sec:prep-vs-proj: through the time axis only, narrowing
%             as the coupling grows with the modes.
% THREAD      Each extension is carried as a coupled linear system on one state,
%             a block per component; the one route that is not, Wu et al.'s
%             doubled space, only projects.
% CRITERION   Formulation, "sufficient detail to permit readers to create
%             equivalent models": the rows, the data and the circuits of a system
%             and of a lift rebuildable from the body plus Appendix E.
% SOURCES     ThesisPlanning.md §5.6, §3.4, §8; pihm/docs/02_methods.md §3.2-3.5,
%             §6.2; pihm/docs/companion/{21,51,55,60,61,62,63,64,65};
%             pihm/docs/16_extensions.md; pihm/docs/17_navier_stokes.md;
%             pihm/docs/03_problems/07_r7.md; claims A-10, B-12, B-18, B-19, B-22,
%             C-09, C-10; src/pihm/circuits/{systems,tensor,standard,sources,
%             assemble}.py, ir.py (matching_terms), problems/fluids.py
%             (real_embedding); Wu et al., p. 11 (the Fourier-basis remark).
% EVIDENCE    Plan T-20's rule (v3.4, StructuralReview.md F1): proofs and exact classical
%             computations only (kernels, gaps, alpha/||.|| and term counts of built
%             couplings). Degrees, preparations and gates per query are §8.4-8.5's (the
%             tensor layout's CNOTs: tab:lift-circuits).
% =============================================================================

Every residual so far has carried one real scalar field, the scope Wu et al.\ set themselves. They set complex-valued states aside, because with a Fourier basis, whose differentiation matrix is imaginary, the ground state of the resulting energy would be hard to find \cite{wu2025pihm}, and they carry a quadratic term in the doubled space of \Cref{sec:pihm-doubled}. The doubled space's kernel is degenerate, so there is no gap above the ground value and \Cref{eq:degree} gives no finite degree; as the degree grows, a filter returns the initial state's projection onto the kernel. A filter needs a unique ground state, which leaves the question \emph{can the solution of a nonlinear equation be prepared as the unique ground state of a Hamiltonian, and does the reformulation's advantage survive there?}

Each extension that prepares is the same object, a coupled linear system on one state with one block per component: several fields directly, a complex field by realification as two real ones, and a polynomial nonlinearity by the Carleman lift, whose components are the field's monomials. The system is built first because the lift that answers the question is an instance of it; the doubled space is the one route that is not. The constructions and proofs in full are in \Cref{app:nonlinear}.

\subsection{Coupled First-Order Systems}
\label{sec:systems-fields}

Consider a first-order system of $C$ fields with a constant coupling matrix $\Acoup$, $z_i'=\sum_j\Acoup_{ij}z_j$. A time interval $t\in[0,T]$ is carried on the Chebyshev axis $\xi=2t/T-1$, so $dz/d\xi=\tfrac T2\Acoup z$ and one ground state holds the whole trajectory, as in the Feynman--Kitaev clock of Barison et al.\ \cite{barison2022variational} and the space--time ground state that Pool et al.\ compute for Burgers' equation \cite{pool2024nonlinear}, both found variationally. Here the time axis is a Chebyshev series and the ground state is prepared by a filter. The state holds one block of $N$ Chebyshev coefficients per field, the field index the most significant, on a register of $\lceil\log_2C\rceil$ qubits; where $C$ is not a power of two, padding blocks decouple, as they do in the chain (\Cref{app:systems}, with R7's rows). In the standard form (\Cref{def:standard-form}) the state holds the fields alone and every derivative is eliminated through the dense differentiation matrix $\G$; the descriptor carries a system in the mass form of \Cref{eq:mass-form}. The equation rows are
\begin{equation}
    \underbrace{I_C\otimes\G_\xi-\tfrac T2\,\Acoup\otimes I_N}_{\text{standard form}},
    \qquad
    \underbrace{I_C\otimes\Dt_\xi-\tfrac T2\,\Acoup\otimes\Bt_\xi}_{\text{descriptor (mass form)}},
    \label{eq:system-rows}
\end{equation}
the second being the first multiplied on the left by $I_C\otimes\Bt_\xi$, so by \Cref{thm:kernel-equivalence} their kernels coincide.

With the data rows of \Cref{sec:systems-data}, an \gls{ivp} has a one-dimensional solution space (\Cref{prop:ivp-uniqueness}). Its solution is entire, so its residual on polynomials of degree below $N_t$ falls spectrally. A truncated lift is a regular first-order system (\Cref{def:regular}), so \Cref{thm:flat-gap} bounds the descriptor residual's second singular value away from zero as $N_t\to\infty$, in either layout and whether or not the truncation converges, but without a constant. At a finite $n_t$, and for the standard form, no bound is proved, and the reference's trust test (\Cref{sec:pipeline-reference}) checks the near-kernel's dimension on every run.

The descriptor's chain could instead carry each field's derivative as a block, as it does for a field of higher order. Such a block doubles the state and, even with the rescaling of \Cref{sec:descriptor-rescaling}, costs gap: on the harmonic oscillator $u'=v$, $v'=-\pi^2u$ (R7) at $n=8$, the mass form's gap is $\res[3]{chain-ratio/build/R7/descriptor/n8}{gap_ratio}$ times the block-rescaled chain's, and between $\res[3]{chain-ratio/build/R7/descriptor/range}{low}$ and $\res[3]{chain-ratio/build/R7/descriptor/range}{high}$ times from $n=3$ upwards. The chain is kept only as a labelled ablation, and every system and lift in the results is compared in the mass form. A system of higher order enters by the usual reduction to first order, as R7 is $u''=-\pi^2u$ written as two fields; none of higher order than R7's was run.

\subsection{Data Encoding for Systems}
\label{sec:systems-data}

A single field's initial or boundary value is imposed by a homogeneous zero row, which leaves the field's scale free, and a regular datum, which fixes it (\Cref{sec:standard-constraints}). Vector data $y=(y_1,\dots,y_C)$ at a point $x^*$ need the same separation. With an anchor component $a$, $y_a\ne0$, and $y_\perp=\sum_{c\ne a}y_c\ket{c}$ the rest of the data, the \emph{ratio row} on the field register is
\begin{equation}
    \Frat = y_a\bigl(I-\ket{a}\bra{a}\bigr)-y_\perp\bra{a},
    \qquad
    (\Frat v)_c = y_a v_c - y_c v_a\ \ (c\ne a),\quad (\Frat v)_a=0,
    \label{eq:ratio-row}
\end{equation}
and the data rows are $(\Frat/y_a)\otimes\Brow(x^*)$. The kernel of $\Frat$ is exactly the line through $y$, so the rows impose $z_c(x^*)=(y_c/y_a)\,z_a(x^*)$: the direction of the data, with their scale free. The regular datum on the anchor fixes the scale, and a component whose datum is zero is held to zero exactly. Divided by $y_a$, the rows' non-zero singular values are $\|y\|/|y_a|$ and, for $C>2$, one, each times the zero row's norm, so the anchor with the largest $|y_a|$ conditions them best (\Cref{prop:ratio-rows}; their three-term circuit is in \Cref{app:systems}). A system's initial state weights each field's block by the stated data $w_0/\|w_0\|$ (\Cref{sec:pipeline-filter}); equal weights would cancel exactly on the Taylor--Green vortex of \Cref{sec:fluids}, whose two modes carry data of opposite sign.

\subsection{Complex-Valued Linear Equations}
\label{sec:systems-complex}
% TRAPS: "handles complex equations" unscoped; "a complex-valued field was prepared"; the tensor
%   layout called "prepared"; the legacy DESolverLib cited; the section billed above its evidence
%   (2026-10-14 review: the real Fourier bases of R9-R10 lead, R7 is the smallest instance).

The Fourier bases that led Wu et al.\ to set complex-valued states aside are carried here on real bases. The fluids of \Cref{sec:fluids} use sine modes for Burgers' odd solution, and $\cos(\mathbf k\cdot\mathbf x)$ and $\sin(\mathbf k\cdot\mathbf x)$ for the vorticity, an isometric real embedding of the complex Fourier variables in which the truncated system is real (\Cref{app:realification}). In these problems $k$ denotes a wavenumber and $\mathbf k$ a wavevector; elsewhere $k$ is an equation's order. With complex coefficients, data or sources the stacked residual $R$ is complex, yet $H=R^\dagger R$ is Hermitian and positive semi-definite with the smallest right singular vector of $R$ as its ground state, and the reflection and filter of \Cref{sec:pipeline} need only $H$. The filter never sees $R$. Every residual prepared here is real, so this describes the formulation and no run.

A complex linear equation is carried instead by realification. Writing $z=u+\mathrm{i}v$, it becomes a real system of twice as many fields; for one mode $z'=cz$ with a complex coefficient $c$,
\begin{equation}
    u'=\operatorname{Re}c\,u-\operatorname{Im}c\,v,
    \qquad
    v'=\operatorname{Im}c\,u+\operatorname{Re}c\,v,
    \label{eq:realification}
\end{equation}
a first-order system with one more field-register qubit. The map $z\mapsto(\operatorname{Re}z,\operatorname{Im}z)$ is a real-linear isomorphism that carries solutions to solutions, and real ratio rows on the realified data fix the phase with the direction, so an \gls{ivp} keeps a one-dimensional solution space (\Cref{app:realification}). R7 is the smallest instance: with $z=\pi u+\mathrm{i}v$, its equations are the real and imaginary parts of $z'=-\mathrm{i}\pi z$, a single-mode Schr\"odinger-type equation. The one natively complex encoding built is the tensor layout of the Carleman coupling (\Cref{sec:ext-circuits}), whose \gls{lcu} weights carry the phases $\mathrm{i}^{\sum_ja_j}$ of the Fourier symbols $\prod_jk_j^{a_j}$ it applies; it encodes the coupling alone and prepares no state. Native complex preparation would also need complex sources, data and coefficients in the circuits' loaders, which are not built, and is left to further work (\Cref{sec:future-work}).

\subsection{Doubled Space}
\label{sec:doubled-space}

Wu et al.'s doubled space is not a linear system on one state. We run it in the standard form only, the form in which they state it; the descriptor's route to a nonlinear equation is the Carleman lift, and the proposition below applies to any Hamiltonian on the doubled space. Its encoding is exact, each product folded to one Chebyshev series by the nodal fold of \Cref{sec:ext-circuits} and each linear term carried by a copy of the collapse row (\Cref{app:doubled-space}), and its kernel is large: the equation block has rank at most $2N$, the number of coefficients of the folded register, and each of the $m$ invariant constraints adds at most one, so
\begin{equation}
    \dim\ker H_\otimes \;\ge\; N^2-2N-m,
    \label{eq:doubled-kernel}
\end{equation}
above Wu et al.'s count of $2^{2n-1}=N^2/2$ from $n=3$ for $m\le2$; the dimensions computed on R6 are in \Cref{tab:doubled-kernel}.

\begin{proposition}[A filter projects onto the doubled kernel]
\label{prop:no-penalty}
Let $H_\otimes\succeq0$ act on $\mathbb{R}^{N^2}$, $N\ge2$, with $\Pi$ the projector onto its kernel.
\begin{enumerate}
    \item The symmetric product states $\{u\otimes u\}$ do not form a linear subspace, and a kernel does. A kernel that contains two independent product states also contains states that are not products, and no Hamiltonian has the product states as its kernel.
    \item Let $P_d$ be filters with $P_d(0)=1$ whose largest magnitude on the rest of the spectrum tends to zero, as for the minimax filter as $d\to\infty$ or for $e^{-\beta H_\otimes}$ as $\beta\to\infty$. For any $\psi_{\mathrm{init}}$ with $\Pi\psi_{\mathrm{init}}\ne0$,
    \begin{equation}
        \lim_{d\to\infty}\frac{P_d(H_\otimes)\,\psi_{\mathrm{init}}}{\bigl\|P_d(H_\otimes)\,\psi_{\mathrm{init}}\bigr\|}
        =\frac{\Pi\,\psi_{\mathrm{init}}}{\bigl\|\Pi\,\psi_{\mathrm{init}}\bigr\|}.
        \label{eq:doubled-projection}
    \end{equation}
    \item A penalty $P\succeq0$ with $\ker(H_\otimes+P)=\operatorname{span}\{\psi_*\otimes\psi_*\}$ must vanish on $\psi_*\otimes\psi_*$ and on no other vector of $\ker H_\otimes$: a condition stated in terms of the solution $\psi_*$.
\end{enumerate}
\end{proposition}
The proof is given in \Cref{app:proof-no-penalty}. A filter returns the initial state's projection onto a kernel that holds the solution's product state among exponentially many others, so $\psi_{\mathrm{init}}$ decides which state comes back, and the condition that a penalty singles out the solution is stated in terms of that solution, which is what we mean by answer-dependent. The proposition does not show that no penalty built from the problem's data alone can meet the condition, and none is built here. Wu et al.'s selection of ``a degenerate state that closely matches the analytical solution'' \cite{wu2025pihm} is such a choice, and the doubled space is accordingly reported here as a projection.

\subsection{Carleman Lift in Space and Time}
\label{sec:carleman}
% TRAPS: re-deriving the space-time rows (they are eq:system-rows with A -> A_K); the lift's
%   preparation stated without the convergence caveat; rho < 1 as necessary.

The Carleman lift of a polynomial system is a linear system on the system's monomials \cite{carleman1932,kowalskisteeb1991,liu2021carleman,krovi2023improved} (\Cref{sec:bg-carleman}), so it is built from the same rows and data rows as any other system: the components are monomials, the ratio rows carry the data's moments at $t=0$, and $\Acoup$ is replaced by the lifted generator $\Acoup_K$ in \Cref{eq:system-rows}, in either form. For $u\in\mathbb{R}^M$,
\begin{equation}
    \frac{du}{dt}=\sum_{p\ge1}F_p\,u^{\otimes p}
    \quad\longrightarrow\quad
    \frac{dz}{dt}=\Acoup_K\,z,
    \qquad
    z=\bigl(u_\alpha\bigr)_{1\le|\alpha|\le K},\quad u_\alpha=\textstyle\prod_iu_i^{\alpha_i},
    \label{eq:carleman-lift}
\end{equation}
where the state holds every monomial of degree $1$ to $K$ once, level by level (the symmetric layout), and each row of $\Acoup_K$ differentiates a monomial and substitutes $u'$, dropping the terms beyond degree $K$ (\Cref{app:carleman-layouts}). The lift is carried in the scaled variable $u/s$, $s=\max(\|u_0\|,1)$, a diagonal similarity of $\Acoup_K$ that keeps the data's moments on level $\ell$ from growing as $\|u_0\|^\ell$ (\Cref{prop:level-scale}).

With $\lambda_1$ the eigenvalue of $F_1$ of largest real part and $\|F_2\|$ taken on symmetric tensors,
\begin{equation}
    \rho=\frac{\|u_0\|\,\|F_2\|}{\bigl|\operatorname{Re}\lambda_1(F_1)\bigr|}
    \label{eq:carleman-ratio}
\end{equation}
is the ratio of Liu et al., who prove for a dissipative quadratic system that $\rho<1$ suffices for the truncation error at order $K$ to fall as $\mathcal{O}(\rho^K)$, uniformly in time \cite{liu2021carleman}. The condition is sufficient and a worst case over all data of norm $\|u_0\|$; nothing in the bound says the lift diverges at $\rho\ge1$, and other analyses give conditions of other forms: explicit time-dependent error bounds \cite{forets2017carleman}, a milder ratio in the $\ell^\infty$ norm for reaction--diffusion equations \cite{liu2023reactiondiffusion}, and exponential convergence over all time for a class of stable systems \cite{amini2025carleman}. For $y'=-y+y^2$ from $y(0)=\tfrac12$ (R8), $\rho=\tfrac12$; on Burgers it grows with the number of modes (\Cref{sec:fluids}). Convergence is measured on every run, order by order against the exact solution of the truncated lift (\Cref{sec:results-nonlinear}).

The truncated lift is a linear \gls{ivp} with its data fixed up to scale, so the same argument gives a near-kernel expected to be one-dimensional in both forms and for every $\Acoup_K$ (\Cref{app:carleman-spacetime}). Either form inherits the uniqueness from the lifted \gls{ivp}. The decoded field's error splits in two,
\begin{equation}
    \|\hat u-u\|_\infty\;\le\;\underbrace{\bigl\|z^{(K)}_1-u\bigr\|_\infty}_{\text{lift truncation}}+\underbrace{\bigl\|\hat u-z^{(K)}_1\bigr\|_\infty}_{\text{time-axis discretisation}},
    \label{eq:carleman-error}
\end{equation}
with $z^{(K)}_1$ the first level of the truncated lift's exact solution. The lift truncation does not depend on $n_t$ and the time-axis discretisation falls spectrally to rounding, so the field's error falls with $n_t$ until it meets the floor set by the lift and then holds; every run records both terms.

% CUT-CANDIDATE (page pass 2026-10-14, about 0.5 pp; the author's call, not applied): keep one body
%   paragraph (matchings encoded at the sum of their largest entries, the same factor in both forms;
%   the tensor layout's 2K - 1 terms; the nodal fold exact at alpha = 2) and move the construction
%   detail and the Ucg/DCT budgeting to App. E (app:matchings, app:systems).
\subsection{Circuits for the Coupling and the Products}
\label{sec:ext-circuits}
% 2026-10-16 (GeminiReviewPlan row 8): one sentence comparing the matchings' alpha <= (2delta-1)||Gamma||_max
%   with sparse access (gilyen2019qsvt, STOC §5.2: alpha = sqrt(s_r s_c) ||A||_max; arXiv Lemma 48); the
%   documented alternative is Pauli words, qualitatively (pihm decision 42's numbers are not in results-v1).

Everything in a system's or a lift's rows except its coupling is built from the bands and rank-one rows of \Cref{sec:descriptor-block-encoding}; the ratio rows are three \gls{lcu} terms.

A system's coupling $\Acoup$, or a lift's $\Acoup_K$ in the symmetric layout, is sparse and irregular. We split it into \emph{matchings}, each with at most one entry in every row and column: a permutation of the field register's basis states followed by a diagonal of rotations on one flag qubit. Each is encoded at its largest entry, so the coupling's subnormalisation is the sum of those entries, and the split is a greedy edge colouring (\Cref{app:matchings}). With $\delta$ the most entries in any row or column of the coupling and $\|\Acoup\|_{\max}$ its largest entry in magnitude, the colouring uses at most $2\delta-1$ matchings, so $\alpha\le(2\delta-1)\|\Acoup\|_{\max}$. This is within a factor of two of the $\delta\|\Acoup\|_{\max}$ of a sparse-access encoding \cite{gilyen2019qsvt}, and it avoids that encoding's oracles, which would have to compute the position of each row's entries from the symmetric layout's irregular monomial index. The expansion in Pauli words that encodes the other field-register matrices would need a term for every word with a non-zero coefficient. The factor is the same in both forms, so the coupling costs the same in both.

On Burgers' lift on sixteen modes at $K=3$ the matchings encode $\Acoup_K$ at $\alpha/\|\Acoup_K\|=\res[3]{lift/R9c-M16/standard/n0/carleman=3/verify=ir}{alpha_over_norm}$. On the Navier--Stokes lifts the ratio is larger, $\res[3]{lift/R10c-M80/standard/n0/carleman=2/verify=ir}{alpha_over_norm}$ on eighty modes at $K=2$ (\Cref{tab:lift-circuits}). The number of matchings grows with $\Acoup_K$'s entries, and so with the number of modes.

For a Fourier--Galerkin equation, whose linear operators are Fourier multipliers and whose products are convolutions, a second layout, the tensor layout, carries level $\ell$ as the whole tensor $u^{\otimes\ell}$ in $\ell$ slots of wavevectors. The coupling is then a sum of terms, each applying $F_p$ to $p$ adjacent slots, $2K-1$ of them for a quadratic field whatever the number of modes (\Cref{prop:tensor-terms}). Its weights carry the complex phases of the Fourier symbols, and its cost per query is compared with that of the matchings in \Cref{sec:results-ns}.

A product of two series is folded back to one by transforming to node values, multiplying pointwise and transforming back. The doubled space's nodal fold does this with the Chebyshev--Gauss transform $\Ucg$, exactly and at $\alpha=2$, within a factor $\res[3]{fold/n2}{alpha_ratio}$ of its norm at $n=2$ and $\res[3]{fold/n12}{alpha_ratio}$ at $n=12$: tight to a constant factor, and not optimal (\Cref{prop:fold-exact}). $\Ucg$ is the discrete cosine transform, for which quantum circuits with polylogarithmically many gates are known \cite{klappenecker2001dct}. The fold serves only the doubled space, which prepares no state, so $\Ucg$ is budgeted as a cited primitive and not built gate by gate. The tensor layout's Fourier fold uses the built \gls{qft}, at $\alpha=N^{1/2}$ per axis, the convolution's own norm.

\subsection{Burgers and Navier--Stokes Equations}
\label{sec:fluids}
% NOTATION: k is the wavenumber in this section only (the equation's order elsewhere); bold k is
%   a wavevector. Stated to the reader once, at the end of sec:systems-complex's realification
%   paragraph.

The lift is tested on two periodic flows whose quadratic terms are convolutions. Truncated to finitely many Fourier modes on a real basis (\Cref{sec:systems-complex}) they are polynomial systems of the form of \Cref{eq:carleman-lift}. Truncating to Fourier modes adds a third error beside the two of \Cref{eq:carleman-error}: the Galerkin truncation, measured against the equation itself, whereas the lift truncation is measured against the Galerkin system's \gls{ivp}. The systems, their couplings and the parameters of every run are given in \Cref{app:fluids}.

\paragraph{Viscous Burgers (R9).} On $x\in[0,2\pi)$,
\begin{equation}
    u_t+u\,u_x=\nu\,u_{xx},\qquad u(x,0)=a\sin x,
    \label{eq:burgers}
\end{equation}
whose solution stays odd and lives on the sine modes $\sin kx$, $k\le M$. The Galerkin system has $F_1=\operatorname{diag}(-\nu k^2)$ and $F_2$ the projection of $-\tfrac12(u^2)_x$. The least damped mode is $k=1$, so $\rho=a\|F_2\|/\nu$, which grows with $M$ through $\|F_2\|$. The runs vary the strength $a/\nu$ and the number of modes $M$, and the Cole--Hopf solution \cite{hopf1950burgers,cole1951burgers} measures the truncation error.

\paragraph{Two-Dimensional Navier--Stokes (R10).} In vorticity form on the periodic square,
\begin{equation}
    \omega_t+J(\psi,\omega)=\nu\,\nabla^2\omega,
    \qquad
    \nabla^2\psi=-\omega\ \ \text{on }[0,2\pi)^2,
    \qquad
    J(\psi,\omega)=\psi_x\omega_y-\psi_y\omega_x,
    \label{eq:navier-stokes}
\end{equation}
with $\nu=1$, the vorticity on the real modes $\cos(\mathbf k\cdot\mathbf x)$ and $\sin(\mathbf k\cdot\mathbf x)$ of the half-plane with $0<|\mathbf k|^2\le k_c^2$, the stream function eliminated as $\hat\psi_{\mathbf k}=\hat\omega_{\mathbf k}/|\mathbf k|^2$, and every coupling rational and exact. Modes of equal $|\mathbf k|$ do not interact, and without viscosity the truncation conserves the energy and the enstrophy (\Cref{app:fluids}); a pseudo-spectral \gls{dns} measures the truncation error. The Taylor--Green vortex $\omega_0=2\sin x\sin y$ \cite{taylorgreen1937} occupies two modes of equal $|\mathbf k|$, so $\omega=\omega_0e^{-2\nu t}$ exactly and its lift is exact at every order (\Cref{prop:exact-lift}): R10a tests the whole chain at a known answer and leaves the lift's convergence untested. A counter-rotating vortex pair of norm $\|w_0\|=a$ spans several $|\mathbf k|$, so its triads are active and its lift is tested.

\subsection{Comparison of the Carleman and Doubled-Space Routes}
\label{sec:prep-vs-proj}

\begin{table}[!htbp]
    \centering
    \caption[Preparation by the Carleman lift against projection by the doubled space]{\textbf{The Carleman lift prepares; the doubled space projects.} Each nonlinear route by the forms it is run in, its kernel, what a filter returns from an initial state $\psi_{\mathrm{init}}$, and what a trusted kernel (one-dimensional, gap ratio above 20) certifies. $N=2^n$; $m$, number of invariant constraints; \gls{ivp}, initial-value problem.}
    \label{tab:prep-vs-proj}
    \small
    \begin{tabular}{@{}p{0.16\linewidth} p{0.12\linewidth} p{0.2\linewidth} p{0.2\linewidth} p{0.2\linewidth}@{}}
        \toprule
        Route & Forms run & Kernel & A filter returns & A trusted kernel certifies \\
        \midrule
        Doubled space (Wu et al.; R6) & standard only & degenerate, dimension $\ge N^2-2N-m$ & the projection of $\psi_{\mathrm{init}}$ onto the kernel & nothing: the kernel is never one-dimensional, so never trusted \\
        Carleman lift (R8--R10) & both & one-dimensional near-kernel, from the lifted \gls{ivp}'s uniqueness & the truncated lift's solution & the lifted \emph{linear} problem, not the lift's convergence \\
        \bottomrule
    \end{tabular}
\end{table}

A trusted kernel certifies that the residual singles out the truncated lift's solution (last column of the table). It says nothing about the distance of that solution from the nonlinear one, the lift truncation error, which is measured separately. The two come apart on Burgers' lift: past its convergence boundary the kernel stays trusted while the lift diverges (\Cref{sec:results-nonlinear}).

On a lift the descriptor's effect enters through the time axis. At one Carleman order the two forms differ only by the left factor $I\otimes\Bt_\xi$ (\Cref{eq:system-rows}), so the time axis contributes $\|\Dt_\xi\|=2(N_t-1)$ to the norm of the mass form against $\|\G_\xi\|=\Theta(N_t^2)$ to that of the standard form. The coupling term has norm $\tfrac T2\|\Acoup_K\|$ in the standard form and $\|\Bt_\xi\|_2\in[2,2+\sqrt2]$ times that in the mass form (\Cref{sec:descriptor-construction}); it is of the same order in both and grows with the modes. The truncated lift is a regular first-order system (\Cref{def:regular}), the class for which \Cref{eq:descriptor-degree} predicts $d=\tilde\Theta(N_t)$ at a fixed truncation, so the mass form's degree should grow more slowly with $N_t$. At fixed $n_t$ the advantage should narrow as the truncation widens and the coupling term comes to dominate both norms, and on a coarse time axis the factor $\|\Bt_\xi\|_2$ can make the mass form's norm the larger. \Cref{sec:results-nonlinear,sec:results-ns} measure both effects, from the scalar R8 to two-dimensional Navier--Stokes.

% CHECKLIST (marker)
%   1. Is the scope gap stated from Wu et al.'s own words, and Q(iii) posed at the head?
%   2. Is the thread (one state, a coupled linear system) visible in the opening and in
%      sec:carleman's first sentence?
%   3. Are systems and lifts rebuildable from body and Appendix E?
%   4. Is the mass form applied, not redefined (sec:descriptor-layouts owns it)?
%   5. Is realification stated as the route run, with R7 as its instance, and every complex-valued
%      claim scoped (formulation; one encoding verified; preparation not built)?
%   6. Is "prepared" never used for the doubled space?
%   7. Is rho < 1 stated as SUFFICIENT, not necessary?
%   8. Is the nodal fold's alpha = 2 called tight (within 1.27 of ||n_1||_2 at n = 2), not optimal?
%   9. Is the DCT named as unbuilt, with the reason?
%  10. Is uniqueness attributed to Carleman (the lifted IVP), not to a form?
%  11. Is trust distinguished from convergence (sec:prep-vs-proj)?
%  12. Is each construction stated once in the chapter (matchings, ratio rows, the time axis)?
% TRAPS (chapter): "handles complex equations" unscoped; any complex-valued result resting on the
%   legacy DESolverLib (never cited); redefining the mass form or restating its kernel theorem;
%   calling the tensor encoding "prepared"; claiming higher-order systems were run;
%   "vector-valued solver"; "K^2 terms" (it is 2K - 1); "both forms" for the doubled space (it is
%   standard-form only); rho < 1 as necessary; "optimal" for the fold; re-deriving the
%   space-time rows.
